\section{Discussion and Limitations}
Layerwise rates affect the metric through which pretraining selects directions. Our certified construction shows that this can produce a fixed rise in an independent task's optimal population projection risk while the loss decreases. The controlled high-dimensional experiments then distinguish two effects: reselecting directions within a fixed encoder space and changing the encoder space itself. Neither effect should be inferred from a projection curve alone.

The population distinction is conditional. Finite data break exact isotropy, empirical nonlinear moment objectives need not have unbiased gradients, and finite steps contribute additional changes. Linear encoder movement under complete empirical training has a reproducible component; describing it as only zero-mean stochastic noise would be incorrect. Likewise, early success of population-based risk predictions does not imply a reliable late-stage stochastic intervention. We report both failures explicitly.

The theorem is computer-assisted rather than a proof-assistant formalization. Its continuous initialization box is narrow, and the high-dimensional random-seed experiments do not enlarge its certified domain. Its radial example is not contrastive-specific: smooth decreasing radial objectives share the same energy-parameterized path over intervals where their scalar speed remains positive. Full-support coordinates are permutation-equivariant and do not by themselves explain a universal low-to-high frequency order. Local-kernel experiments add a support constraint, not a new locality theorem.

Natural-image evidence is deliberately limited. The shallow grayscale architecture, finite training budget, and small seed counts do not support general claims about visual backbones. Fixed-encoder rank-8 fitting shows that head selection can discard useful class directions, but it uses labels and the gap already exists at initialization. Covariance effective-rank concentration is also not equivalent to loss of all linearly readable information. The current evidence supports a controlled account of task-dependent representation selection and a measurable boundary on its predictability.

\section{Conclusion}
Actual factorized training can improve its pretraining objective while first improving and then worsening the optimal population risk of a compressed representation. We certify this behavior on an explicit initialization box and identify a matched-function learning-rate threshold for its initial direction. Population and empirical experiments separate projection selection from encoder subspace motion, while exposing finite-data drift and late-stage prediction failure. CIFAR-10 diagnoses substantial head-selection gaps without demonstrating a stable unsupervised ratio benefit. These distinctions make clear which parts of the phenomenon are explained, which depend on idealized assumptions, and which remain unresolved.
